\section{Introduction: the idea of a layer}
The SWAFL application (Soil Water, Atmosphere, Feedback and Layout) is organised in \emph{layers}. A layer is a self-contained view of one question about a field: when to irrigate, what a mulch does, how the structure of the soil changes drainage, how a plant grows, how the water balance of the field closes. All layers read the same inputs (the soil, the crop, the weather, the management) and write to a shared summary, so that a change in one place reaches all of them. The layers belong to two groups: those that \emph{read} the archived results of the ten-thousand Daisy simulations (control, mulch, structure, loops) and those that \emph{compute} with an engine (the field balance, the seasonal data, the crops, the designs, the phenotyping). The application also contains the \emph{Daisy workbench}, in which a user plans and reads a Daisy experiment. This chapter tours the layers in the order of the project tree; each section gives the background of the question, the main relation the layer uses, and one number reproduced by hand. The following chapters take each family to depth.

\section{Irrigation design: surface, drip and sprinkler}
The design layer takes a soil, a crop and a field and returns a system that can irrigate it: for surface irrigation the inflow, the advance and the cut-off time; for drip the emitter, the lateral and the wetted bulb; for sprinklers the nozzle, spacing and pressure; for all of them the main pipe and the pump. Its first relation is the \kw{system capacity}: to apply a net depth $d$ over an area $A$ in $T$ days of $h$ working hours each, with application efficiency $E_a$,

\begin{equation}Q=\frac{10\,A\,d}{E_a\,T\,h}\quad(\mathrm{m^3\,h^{-1}};\ A\ \text{in ha},\ d\ \text{in mm}).\label{eq:cap7}\end{equation}
\par\noindent where $Q$ is the system capacity (m$^3$\,h$^{-1}$), $A$ the area irrigated (ha), $d$ the gross depth per irrigation (mm), $E_a$ the application efficiency, $T$ the irrigation interval (d) and $h$ the hours of operation per day.



\exbox{System capacity}{A 40-ha maize field must receive a net 45 mm in 10 days, working 20 h per day, with sprinklers of $E_a=0.75$. Find the system capacity in m$^3$\,h$^{-1}$ and in L\,s$^{-1}$.}\sol \givenline{$E_a=0.75$}\ans \[Q=10\times40\times45/(0.75\times10\times20)=120\] m$^3$\,h$^{-1}$ $=33.3$ L\,s$^{-1}$ \ans That is the flow to be carried by the mainline and delivered by the pump (Chapter~7); at 60 L\,s$^{-1}$ per block it needs 0.6 blocks working at once. \finans

\section{Irrigation control: sensor and trigger}
The control layer plays a season of soil water and shows when a sensor at a chosen depth passes the trigger. The soil column is a daily model; the sensor reads the suction from the retention curve; the trigger is a suction or a depletion. The numbers that come out are the days beyond the trigger, the irrigation applied and the yield, taken for each soil and trigger from the archived experiments (Chapter~4). A simple version of the same logic can be run by hand with a depletion model.

\exbox{Count of irrigations}{A crop uses 5 mm\,d$^{-1}$ for 100 days from a root zone with RAW = 30 mm refilled at each irrigation. How many irrigations are triggered, and what is the net water applied?}\sol The interval is $30/5=6$ d; irrigations on days 6, 12, \dots, 96: $\lfloor100/6\rfloor=16$ irrigations \ans Net water $=16\times30=480$ mm \ans The remaining 4 days draw the store down by 20 mm, which is the end-of-season depletion. \finans

\section{Mulch and evaporation}
A mulch lowers the vapour flux from the soil by the \emph{strength} $s$ (the fraction by which the flux is reduced) and holds a depth of rain, the \emph{capacity} $c$, that evaporates from the mulch itself. The layer reads the Daisy experiments over strength, rain and capacity and shows the evaporation, drainage, irrigation and yield of the mulched run against the bare one, and the fate of the water saved (Chapter~10). To first order, the evaporation from a mulched soil is $E_m\approx(1-s)E_s+E_{int}$, with $E_{int}=\sum\min(P_j,c)$ over the rain events.

\exbox{Evaporation from a mulched soil}{Bare-soil evaporation is 180 mm per season. A mulch with strength 0.6 and capacity 2 mm meets 35 rain events of more than 2 mm. Estimate the evaporation of the mulched soil and the apparent saving.}\sol \[E_m=(1-0.6)\times180+35\times2=72+70=142\ \text{mm}\]
apparent saving $180-142=38$ mm \ans The interception of the mulch gives back 65 per cent of the saving of 108 mm, which is why a mulch of large capacity saves less than its strength suggests. \finans

\section{Soil structure and tillage}
Tillage loosens the topsoil (lowers the bulk density, raises the porosity and the conductivity) and the soil then consolidates with time and rain. The layer shows the effect of tillage (conventional, reduced, none) and of the consolidation rate on drainage, evaporation and yield (Chapter~10). A simple model: the bulk density relaxes exponentially, $\rho_b(t)=\rho_\infty-(\rho_\infty-\rho_0)e^{-kt}$, and the saturated conductivity scales with the Kozeny--Carman factor $K_s\propto\phi^3/(1-\phi)^2$.

\exbox{Consolidation after tillage}{After ploughing, $\rho_b=1.15$ Mg\,m$^{-3}$ and it relaxes towards $1.45$ with $k=0.02$ d$^{-1}$. Find $\rho_b$ at 0, 30, 60, 120 and 240 days and the relative saturated conductivity $K_s(t)/K_s(0)$ ($\rho_s=2.65$).}\sol \givenline{$\rho_b=1.15$ Mg\,m$^{-3}$ ; $k=0.02$ d$^{-1}$ ; $\rho_s=2.65$}\ans $t=0$: $\rho_b=1.150$, $K_s/K_{s0}=1.000$; $t=30$: $\rho_b=1.285$, $K_s/K_{s0}=0.603$; $t=60$: $\rho_b=1.360$, $K_s/K_{s0}=0.455$; $t=120$: $\rho_b=1.423$, $K_s/K_{s0}=0.358$; $t=240$: $\rho_b=1.448$, $K_s/K_{s0}=0.325$ \ans The conductivity of the tilled layer falls to about a third within the first two months, which is the reason the benefit of tillage on drainage is short. \finans

\section{Phenotyping: the deterministic plant and the volumetric scan}
The two phenotyping layers grow a plant by a deterministic crop model and then scan it with a virtual camera rig to measure its traits (Chapter~11). The traits are height, crown size, volume, leaf area and biomass. A convenient geometric model of a plant is an ellipsoid of semi-axes $a$, $a$ (horizontal) and $b$ (vertical): volume $V=\frac43\pi a^2b$, and the leaf area index of the ground area $A_g=\pi a^2$ is $\mathrm{LAI}=\mathrm{LA}/A_g$ with LA the leaf area.

\exbox{Volume and LAI of a plant}{A plant of height $H=0.9$ m has a crown of horizontal radius 0.35 m and is modelled as an ellipsoid. Find its volume and, if its leaf area is 0.85 m$^2$, the LAI of the ground area it covers.}\sol \givenline{$H=0.9$ m has a crown}\ans Semi-axes $a=0.35$ m, $b=H/2=0.45$ m: \[V=\frac43\pi\times0.35^2\times0.45=0.231\ \text{m$^3$}\] \ans Ground area $A_g=\pi a^2=0.385$ m$^2$, so \[\mathrm{LAI}=0.85/0.385=2.21\quad\text{\finans}\]

\section{Field water balance}
The field-balance layer runs the daily bucket (Chapter~9): the depletion $D_r$ of the root zone is updated each day, $D_{r,i}=D_{r,i-1}-P_i-I_i+ET_{a,i}+D_{p,i}$, with deep percolation $D_{p}=\max(0,P+I-ET_a-D_{r,i-1})$ when the soil is refilled beyond field capacity, and $ET_a=K_sK_cET_0$.

\begin{table}[!htbp]\centering\small\caption{One week of the daily bucket: TAW = 120 mm, RAW = 60 mm, $K_c=1.1$, initial depletion 40 mm; all in mm.}
\begin{tabular}{rrrrrr}\toprule
\hd{Day} & \hd{$ET_0$} & \hd{Rain} & \hd{$K_s$} & \hd{$ET_a$} & \hd{$D_r$ end}\\\midrule
1 & 5 & 0 & 1.00 & 5.50 & 45.5\\
2 & 5.5 & 0 & 1.00 & 6.05 & 51.5\\
3 & 6 & 0 & 1.00 & 6.60 & 58.1\\
4 & 6 & 12 & 1.00 & 6.60 & 52.8\\
5 & 5.5 & 0 & 1.00 & 6.05 & 58.8\\
6 & 5 & 0 & 1.00 & 5.50 & 64.3\\
7 & 4.5 & 0 & 0.93 & 4.60 & 68.9\\\bottomrule\end{tabular}\end{table}

Day by day the depletion grows by the actual evapotranspiration; the rain of 12 mm on day 4 lowers it; and as $D_r$ approaches RAW the stress factor $K_s$ starts to fall below 1. Seven days after the start the depletion is 68.9 mm: the crop is beyond the readily available water.

\section{Seasonal data and projections}
The seasonal layers summarise a year of weather and the water balance as charts and tables, and project them forward by a rate of warming and of rainfall change per decade. A projection to the year $y$ after the record applies $\Delta T=r_T\,y/10$ and $P\to P(1+r_P\,y/10)$; the reference evapotranspiration rises by about 2 to 4 per cent per $^\circ$C of warming in humid climates. The aridity index $AI=P/ET_0$ classifies the climate: below 0.2 is arid, 0.2--0.5 semi-arid, 0.5--0.65 dry sub-humid, above 0.65 humid.

\exbox{Aridity index and projection}{A site has $P=560$ mm and $ET_0=1\,120$ mm. The projection is $+0.3$ $^\circ$C and $-2$ per cent rainfall per decade. Find the aridity index now and after 30 years if $ET_0$ rises 3 per cent per $^\circ$C.}\sol \givenline{$P=560$ mm ; $ET_0=1\,120$ mm}\ans Now: $AI=560/1120=0.50$ (dry sub-humid, at the limit with semi-arid). After 30 years: $\Delta T=0.9$ $^\circ$C; \[ET_0=1120\times(1+0.03\times0.9)=1150\ \text{mm}\]
\[P=560\times(1-0.02\times3)=526\ \text{mm}\]
$AI=0.458$ \ans The climate moves to the semi-arid class, and the net irrigation requirement rises by about 8 per cent for a crop of $ET_c\approx0.9ET_0$. \finans

\section{Crops and the crop water requirement}
The crop layer holds the library of 48 crops (Chapter~6) and computes, for a crop, a soil and a climate, the seasonal $ET_c$, the net and gross irrigation requirement, the water productivity and the water footprint. The water productivity is $WP=Y/ET_a$ (kg of yield per m$^3$ of water evapotranspired); typical values are 1.2--2.0 kg\,m$^{-3}$ for maize grain and 0.8--1.2 for wheat.

\exbox{Water productivity}{Maize yields 10.8 Mg\,ha$^{-1}$ with $ET_a=560$ mm; wheat yields 6.4 Mg\,ha$^{-1}$ with $ET_a=420$ mm. Compare their water productivities.}\sol \givenline{$ET_a=560$ mm; wheat yields ; $ET_a=420$ mm}\ans Maize: $WP=10\,800/5\,600=1.93$ kg\,m$^{-3}$; wheat: $6\,400/4\,200=1.52$ kg\,m$^{-3}$ \ans Maize gets more grain per cubic metre of water, as a C4 crop does. \finans

\section{Feedback loops}
The loops layer draws the links between climate, management, soil water, crop and nitrogen as arrows, and runs the Daisy results around them: more irrigation raises the soil water, which raises transpiration and yield but also drainage and nitrogen leaching; a mulch lowers evaporation, which raises the store and so drainage. A loop has a sign (positive if it reinforces, negative if it balances) equal to the product of the signs of its arrows. The nitrogen lost by leaching is $L_N=D\,c_N$ with $D$ the drainage and $c_N$ the concentration of nitrate in the drainage water.

\exbox{Nitrogen lost with drainage}{An extra 120 mm of irrigation raises drainage by 30 mm. The drainage water carries 18 mg\,L$^{-1}$ of nitrate-nitrogen. Find the extra nitrogen lost per hectare and the sign of the loop irrigation $\to$ drainage $\to$ nitrogen loss $\to$ yield.}\sol \[L_N=30\,\mathrm{mm}\times18\,\mathrm{mg\,L^{-1}}=30\,\mathrm{L\,m^{-2}}\times0.018\,\mathrm{g\,L^{-1}}=0.54\] g\,m$^{-2}=5.4$ kg\,ha$^{-1}$ \ans Signs: irrigation $\to$ drainage $(+)$; drainage $\to$ nitrogen loss $(+)$; nitrogen loss $\to$ yield $(-)$; the product is negative, so the loop is a \emph{balancing} loop: the extra water that the crop does not use removes the nitrogen that it needs. \finans

\section{The Daisy model in the application}
The last branch of the project tree is the Daisy workbench, with which a user chooses a Daisy experiment from a ready list or builds one, runs the combinations, reads the results in tables and charts and exports them (Chapters~3 and~9). The experiments of this book are read from the archive of 10\,190 runs; the workbench itself runs on any Daisy installation. The workbench teaches the \emph{logic of an experiment}: a base case, a set of factors with levels, the full factorial of their combinations, and a response. For $k$ factors with $L_i$ levels the number of runs is $\prod L_i$; main effects are the differences of the mean response between levels, and interactions are the departures from additivity.

\exbox{}{A design has three factors: soil (3 levels), mulch strength (3 levels) and rain scale (2 levels). Find the number of runs. In the table of mean yield (Mg\,ha$^{-1}$): mulch 0 = 6.1, mulch 0.5 = 6.5, mulch 1.0 = 6.8; find the main effect of the strongest mulch against none and its relative size.}\sol Runs $=3\times3\times2=18$ \ans Main effect $=6.8-6.1=0.7$ Mg\,ha$^{-1}$ \ans Relative $=0.7/6.1=11.5$ per cent \finans

